% chapters/09/section-02.tex
\section{Menggambar Diagram Keadaan--Transisi}
\label{sec:diagram-keadaan}

Sebelum mesin keadaan ditulis sebagai kode, ia digambar terlebih dahulu.
Gambar itu bukan hiasan: ia adalah rancangan, dan kesalahan pada gambar akan
menjadi kesalahan pada program.

\subsection{Tiga Lambang yang Dipakai}

Diagram keadaan--transisi hanya memakai tiga lambang.

\begin{itemize}
  \item \textbf{Kotak bersudut tumpul} menyatakan sebuah keadaan. Isinya nama
        keadaan itu.
  \item \textbf{Anak panah} menyatakan perpindahan yang mungkin terjadi, dari
        keadaan asal menuju keadaan tujuan.
  \item \textbf{Keterangan pada anak panah} menyatakan syarat perpindahan itu.
        Syaratnya harus dapat diperiksa oleh program, misalnya ``jarak kurang
        dari jari-jari persepsi''.
\end{itemize}

Diagram untuk NPC pada bab ini diperlihatkan pada
Gambar~\ref{fig:diagram-keadaan}. Seluruh perpindahan yang
mungkin tergambar di sana, dan tidak ada perpindahan lain yang sah.

\begin{figure}[htbp]
\centering
% chapters/09/figures/diagram-keadaan.tikz
% Diagram keadaan-transisi NPC: empat keadaan dan tujuh perpindahan yang sah.
% Keterangan pada anak panah dibuat singkat; syarat lengkapnya ada pada tabel
% di dalam teks. Keterangan disusun bertingkat agar tidak saling bertumpuk.
\begin{tikzpicture}[
  keadaan/.style={
    rectangle, rounded corners=4pt, draw=ubPrimary, thick, fill=ubAlt,
    align=center, text width=24mm, inner sep=4pt, minimum height=12mm,
    font=\footnotesize\bfseries
  },
  serang/.style={keadaan, draw=ubRed, fill=ubRed!12},
  panah/.style={-{Stealth[length=2.2mm]}, thick, draw=ubAccent},
  panahbalik/.style={-{Stealth[length=2.2mm]}, thick, draw=ubWarn},
  label/.style={font=\scriptsize, align=center},
  labelbalik/.style={font=\scriptsize, align=center, text=ubWarn}
]

\node[keadaan] (diam) at (0,0) {berdiam};
\node[keadaan] (patrol) at (3.6,0) {berpatroli};
\node[keadaan] (kejar) at (7.2,0) {mengejar};
\node[serang] (serang) at (10.8,0) {menyerang};

% Perpindahan maju, keterangannya disusun bertingkat.
\draw[panah] (diam) -- (patrol);
\node[label, anchor=south] at (1.8,0.78) {1,5 detik};

\draw[panah] (patrol) -- (kejar);
\node[label, anchor=south] at (5.4,1.75) {pemain terlihat};

\draw[panah] (kejar) -- (serang);
\node[label, anchor=south, text width=34mm] at (9.0,0.78)
  {dalam jangkauan serang};

% Perpindahan kembali.
\draw[panahbalik] (serang.south) .. controls (10.8,-1.0) and (7.2,-1.0) ..
  (kejar.south);
\node[labelbalik, anchor=north, text width=34mm] at (9.0,-1.15)
  {keluar dari jangkauan serang};

\draw[panahbalik] (kejar.south) .. controls (7.2,-2.4) and (0,-2.4) ..
  (diam.south);
\node[labelbalik, anchor=north] at (3.6,-2.55) {hilang 1 detik};

\draw[panahbalik] (patrol.north) .. controls (3.6,2.4) and (0,2.4) ..
  (diam.north);
\node[labelbalik, anchor=south] at (1.8,2.55) {putaran selesai};
\end{tikzpicture}

\caption{Diagram keadaan--transisi NPC dengan empat keadaan}
\label{fig:diagram-keadaan}
\end{figure}

\subsection{Tujuh Perpindahan yang Sah}

Diagram tersebut memuat tujuh perpindahan. Tabel~\ref{tab:tujuh-perpindahan}
memuatnya beserta syarat masing-masing, sehingga dapat diperiksa satu per satu.

\begin{table}[htbp]
\centering
\caption{Tujuh perpindahan keadaan beserta syaratnya}
\label{tab:tujuh-perpindahan}
{\setlength{\tabcolsep}{4pt}\footnotesize%
\begin{tabularx}{\textwidth}{@{}C{1.6cm} C{1.6cm} L{3.0cm} Y@{}}
\toprule
\headrow \textbf{Dari} & \textbf{Ke} & \textbf{Jenis syarat} & \textbf{Syarat} \\
\midrule
berdiam & berpatroli & pewaktu & Telah berdiam selama 1,5 detik. \\
\addlinespace
berdiam & mengejar & jarak & Pemain berada di dalam jari-jari persepsi. \\
\addlinespace
berpatroli & berdiam & peristiwa & Satu putaran patroli telah selesai. \\
\addlinespace
berpatroli & mengejar & jarak & Pemain berada di dalam jari-jari persepsi. \\
\addlinespace
mengejar & menyerang & jarak & Pemain berada di dalam jangkauan serang. \\
\addlinespace
mengejar & berdiam & pewaktu & Pemain tidak terlihat selama 1 detik. \\
\addlinespace
menyerang & mengejar & jarak & Pemain keluar dari jangkauan serang. \\
\bottomrule
\end{tabularx}}
\tabsource{Dua keadaan dapat berpindah ke keadaan yang sama
melalui syarat yang sama, misalnya dari berdiam maupun dari berpatroli menuju
mengejar. Hal itu wajar dan tidak menimbulkan persoalan.}
\end{table}

\subsection{Kesalahan yang Sering Terjadi pada Diagram}

Tiga kesalahan berikut paling sering muncul, dan semuanya berpangkal pada satu
sebab yang sama: syarat perpindahan yang tidak lengkap.

\begin{enumerate}
  \item \textbf{Perpindahan tanpa syarat yang jelas.} Menuliskan ``kadang
        beralih mengejar'' tidak dapat diterjemahkan menjadi kode. Setiap
        perpindahan harus memiliki syarat yang dapat diperiksa, yaitu
        perbandingan nilai, keadaan tombol, atau pewaktu.
  \item \textbf{Keadaan tanpa jalan keluar.} Keadaan yang tidak memiliki
        perpindahan keluar membuat NPC terperangkap selamanya. Pada diagram di
        atas, setiap keadaan memiliki sekurang-kurangnya satu perpindahan
        keluar.
  \item \textbf{Syarat yang saling bertentangan.} Apabila dua perpindahan dari
        satu keadaan memiliki syarat yang keduanya benar pada saat yang sama,
        hasilnya bergantung pada urutan pemeriksaan di dalam kode. Gejalanya
        khas: NPC sesekali berperilaku tidak seperti dugaan, dan hanya pada
        keadaan tertentu.
\end{enumerate}

\begin{tipsbox}
Kebiasaan yang berguna: periksa diagram Anda dengan pertanyaan ``apabila
kondisi ini benar, perpindahan mana yang terjadi?'' untuk setiap keadaan dan
setiap nilai jarak yang mungkin. Apabila jawabannya lebih dari satu, urutkan
syaratnya secara sadar dan tuliskan urutan itu pada diagram, misalnya dengan
menomori anak panahnya.
\end{tipsbox}

\subsection{Pewaktu sebagai Syarat Perpindahan}

Tiga dari tujuh perpindahan pada tabel di atas bersyarat
pewaktu, bukan jarak. Pewaktu itu ada tiga jenis, dan ketiganya memiliki
keperluan yang berbeda.

\begin{description}
  \item[Pewaktu berdiam] Menghitung lama NPC berdiam, agar ia tidak langsung
        berpatroli pada bingkai berikutnya. Tanpa pewaktu ini, keadaan berdiam
        tidak memiliki makna karena langsung ditinggalkan.
  \item[Pewaktu tenggang] Menghitung lama NPC telah kehilangan pemain, agar ia
        tidak langsung menyerah. Tanpa pewaktu ini, NPC akan berhenti mengejar
        begitu saja ketika pemain berlari di balik sebuah rintangan selama satu
        bingkai.
  \item[Pewaktu jeda] Menghitung jarak waktu antara dua serangan, agar serangan
        tidak terjadi pada setiap bingkai. Pewaktu ini disebut
        \istilah{cooldown} pada RPS.
\end{description}

Ketiga pewaktu itu dijalankan dengan cara yang sama: ditambah dengan selang
waktu setiap bingkai, lalu dibandingkan dengan nilai batasnya. Cara yang sama
juga dipakai pada pewaktu lain mana pun di dalam permainan.
